\documentclass[12pt,leqno]{article}

\usepackage{amsmath}
\usepackage{hyperref}
\usepackage{logicproof}

\setlength{\textheight}{9.1in}
\setlength{\topmargin}{-.1in}
\setlength{\headsep}{0in}
\setlength{\oddsidemargin}{-.1in}
\setlength{\textwidth}{6.5in}

\setlength{\parskip}{1em}
\setlength{\parindent}{0pt}

% for logicproof
\setlength{\subproofhorizspace}{.5em}
\setlength{\intersubproofvertspace}{.5em}

% proof rules
\newcommand{\Intro}[1]{{#1}{\text{I}}}
\newcommand{\Elim}[1]{{#1}{\text{E}}}
\newcommand{\ElimA}[1]{{#1}{\text{E}_1}}
\newcommand{\ElimB}[1]{{#1}{\text{E}_2}}

% truth tables
\newcommand{\T}{\mathtt{T}}
\newcommand{\F}{\mathtt{F}}

\title{CS 511 Homework \#1}
\author{Anthony DeRossi}
\date{12 Sep 2024}

\begin{document}
\maketitle

\section*{Exercise 1}

[LCS, page 79]: Exercise 1.2.1, parts (h), (i), (j)

Prove the validity of the following sequents:

\begin{itemize}

\item[(h)] $p \vdash (p \to q) \to q$

\begin{logicproof}{1}
  p                & premise \\
  \begin{subproof}
    p \to q        & assumption \\
    q              & $\Elim{\to}$ 2, 1
  \end{subproof}
  (p \to q) \to q  & $\Intro{\to}$ 2--3
\end{logicproof}

\item[(i)] $(p \to r) \land (q \to r) \vdash p \land q \to r$

\begin{logicproof}{2}
  (p \to r) \land (q \to r)  & premise \\
  p \to r                    & $\ElimA{\land}$ 1 \\
  \begin{subproof}
    p \land q                & assumption \\
    p                        & $\ElimA{\land}$ 3 \\
    r                        & $\Elim{\to}$ 2, 4
  \end{subproof}
  p \land q \to r            & $\Intro{\to}$ 3--5
\end{logicproof}

\newpage
\item[(j)] $q \to r \vdash (p \to q) \to (p \to r)$

\begin{logicproof}{3}
  q \to r                  & premise \\
  \begin{subproof}
    p \to q                & assumption \\
    \begin{subproof}
      p                    & assumption \\
      q                    & $\Elim{\to}$ 2, 3 \\
      r                    & $\Elim{\to}$ 1, 4
    \end{subproof}
    p \to r                & $\Intro{\to}$ 3--5
  \end{subproof}
  (p \to q) \to (p \to r)  & $\Intro{\to}$ 2--6
\end{logicproof}

\end{itemize}

\section*{Exercise 2}

[LCS, page 84]: Exercise 1.4.2, parts (g), (h), (i)

Compute the truth table of the formula

\begin{itemize}

\item[(g)] $((p \to q) \to p) \to p$

\[
\begin{array}{ c | c || c | c | c }
  p & q & p \to q & (p \to q) \to p & ((p \to q) \to p) \to p \\
\hline
  \T & \T & \T & \T & \T \\
\hline
  \T & \F & \F & \T & \T \\
\hline
  \F & \T & \T & \F & \T \\
\hline
  \F & \F & \T & \F & \T \\
\end{array}
\]
\medskip

\item[(h)] $((p \lor q) \to r) \to ((p \to r) \lor (q \to r))$

\[
\begin{array}{ c | c | c || c | c | c | c | c | c }
  & & & & & & & & ((p \lor q) \to r) \to \\
  p & q & r & p \lor q & p \to r & q \to r & (p \lor q) \to r &
    (p \to r) \lor (q \to r) & ((p \to r) \lor (q \to r)) \\
\hline
  \T & \T & \T & \T & \T & \T & \T & \T & \T \\
\hline
  \T & \T & \F & \T & \F & \F & \F & \F & \T \\
\hline
  \T & \F & \T & \T & \T & \T & \T & \T & \T \\
\hline
  \T & \F & \F & \T & \F & \T & \F & \T & \T \\
\hline
  \F & \T & \T & \T & \T & \T & \T & \T & \T \\
\hline
  \F & \T & \F & \T & \T & \F & \F & \T & \T \\
\hline
  \F & \F & \T & \F & \T & \T & \T & \T & \T \\
\hline
  \F & \F & \F & \F & \T & \T & \T & \T & \T \\
\end{array}
\]
\medskip

\newpage
\item[(i)] $(p \to q) \to (\lnot p \to \lnot q)$

\[
\begin{array}{ c | c || c | c | c | c | c }
  p & q & \lnot p & \lnot q & p \to q & \lnot p \to \lnot q &
    (p \to q) \to (\lnot p \to \lnot q) \\
\hline
  \T & \T & \F & \F & \T & \T & \T \\
\hline
  \T & \F & \F & \T & \F & \T & \T \\
\hline
  \F & \T & \T & \F & \T & \F & \F \\
\hline
  \F & \F & \T & \T & \T & \T & \T \\
\end{array}
\]
\medskip

\end{itemize}

\section*{Problem 1}

[LCS, page 87]: Exercise 1.5.3, parts (b), (c)

An adequate set of connectives for propositional logic is a set such that for
every formula of propositional logic there is an equivalent formula with only
connectives from that set. For example, the set $\{\lnot, \lor\}$ is adequate
for propositional logic, because any occurrence of $\land$ and $\to$ can be
removed by using the equivalences $\phi \to \psi \equiv \lnot \phi \lor \psi$
and $\phi \land \psi \equiv \lnot (\lnot \phi \lor \lnot \psi)$.

\begin{itemize}

\item[(b)] Show that, if $C \subseteq \{\lnot, \land, \lor, \to, \bot\}$ is
adequate for propositional logic, then $\lnot \in C$ or $\bot \in C$. (Hint:
suppose $C$ contains neither $\lnot$ nor $\bot$ and consider the truth value
of a formula $\phi$, formed by using only the connectives in $C$, for a
valuation in which every atom is assigned $\T$.)

\medskip
Suppose $C \equiv \{\land, \lor, \to\}$. For $C$ to form an adequate set of
connectives it is enough to show that $\lnot$ and $\bot$ can be defined using
the connectives in $C$.

However, we can see from the truth table that a formula $\phi$ formed using
only the connectives in $C$ cannot be used to define $\lnot$.

\[
\begin{array}{ c || c | c | c }
  \phi & \phi \land \phi & \phi \lor \phi & \phi \to \phi \\
\hline
  \T & \T & \T & \T \\
\hline
  \F & \F & \F & \T \\
\end{array}
\]
\medskip

When $\phi$ is $\T$ it is not possible to write a formula logically equivalent
to $\lnot \phi$, since all combinations of connectives are $\T$.

\newpage
\item[(c)] Is $\{\leftrightarrow, \lnot\}$ adequate? Prove your answer.

\medskip
No.

We can prove by induction that a formula using these connectives cannot be
logically equivalent to the formula $\phi \land \psi$ using the property that
all truth tables contain the same number of $\T$ and $\F$ entries.

\begin{center}
\begin{minipage}{.4\linewidth}
  \[
  \begin{array}{ c | c || c | c }
    \phi & \psi & \lnot \phi & \phi \leftrightarrow \psi \\
  \hline
    \T & \T & \F & \T \\
  \hline
    \T & \F & \F & \F \\
  \hline
    \F & \T & \T & \F \\
  \hline
    \F & \F & \T & \T \\
  \end{array}
  \]
  \medskip
\end{minipage}
\quad
\begin{minipage}{.4\linewidth}
  \[
  \begin{array}{ c | c || c }
    \phi & \psi & \phi \land \psi \\
  \hline
    \T & \T & \T \\
  \hline
    \T & \F & \F \\
  \hline
    \F & \T & \F \\
  \hline
    \F & \F & \F \\
  \end{array}
  \]
  \medskip
\end{minipage}
\end{center}

For the base case we observe that every formula of two propositional variables
has the same number of $\T$ and $\F$ entries.

Assume a formula $\phi_i$ is represented by a column of the truth table. We
can extend this formula by one connective to produce a formula $\phi_{i+1}$
which is also represented by a column of the truth table. By induction the
property holds for every formula.

The truth table for $\phi \land \psi$ contains a single $\T$, therefore it is
not possible to construct a logically equivalent formula using only the
connectives $\leftrightarrow$ and $\lnot$.

\end{itemize}

\vspace{3\baselineskip}
\hrule

Solutions to the Lean 4 exercises are available on GitHub:

\url{https://github.com/anthonyde/math2001/blob/cs511-fall2024/Math2001/CS511_Homework/hw1.lean}

\section*{Exercise 3}

\section*{Exercise 4}

\section*{Problem 2}

\end{document}
